% !TEX root = ../../main.tex
%% Corrigé du défi du chapitre 3 — vérifié par verifications/defis_01_05.py.

\begin{blocLivre}{cor-3-25a}
\begin{corrige}{25}{Le théorème de Napoléon}
\demarche{Un triangle équilatéral direct s'écrit $a+jb+j^{2}c=0$ : toute la
géométrie du problème devient du calcul dans $\Cx$, avec la seule règle
$j^{3}=1$, $1+j+j^{2}=0$.}
\partieExo{A}{Une caractérisation des triangles équilatéraux}
\begin{questions}
  \item $j^{3}=\mathrm{e}^{2\mathrm{i}\pi}=1$ ; comme $j\neq1$,
    $1+j+j^{2}=\dfrac{1-j^{3}}{1-j}=0$ ; et
    $\overline{j}=\mathrm{e}^{-2\mathrm{i}\pi/3}=\mathrm{e}^{4\mathrm{i}\pi/3}=j^{2}$.
  \item $(a,b,c)$ est équilatéral direct si et seulement si la rotation
    de centre $a$ et d'angle $\frac{\pi}{3}$ envoie $b$ sur $c$, soit
    $c-a=\mathrm{e}^{\mathrm{i}\pi/3}(b-a)$. Or
    $\mathrm{e}^{\mathrm{i}\pi/3}=-j^{2}$ ($j^{2}=\mathrm{e}^{4\mathrm{i}\pi/3}
    =-\mathrm{e}^{\mathrm{i}\pi/3}$), donc la condition s'écrit
    $c=(1+j^{2})a-j^{2}b=-ja-j^{2}b$. En multipliant par $j^{2}$ (et
    $j^{4}=j$) : $j^{2}c=-a-jb$, soit $\res{a+jb+j^{2}c=0}$.
\end{questions}
\partieExo{B}{Le théorème}
\begin{questions}[resume]
  \item $A'$ est l'image de $C$ par la rotation de centre $B$ et d'angle
    $-\frac{\pi}{3}$ (vers l'extérieur, puisque $ABC$ est direct) ;
    réciproquement, $C$ est l'image de $A'$ par la rotation de centre $B$
    et d'angle $\frac{\pi}{3}$ : $c-b=\mathrm{e}^{\mathrm{i}\pi/3}(a'-b)$.
    Ainsi $(b,a',c)$ est équilatéral direct, $b+ja'+j^{2}c=0$, d'où
    \[
    \res{a'=-j^{2}b-jc},\qquad b'=-j^{2}c-ja,\qquad c'=-j^{2}a-jb .
    \]
  \item $3p=b+c+a'=(1-j^{2})b+(1-j)c$, et de même
    $3q=(1-j^{2})c+(1-j)a$, $3r=(1-j^{2})a+(1-j)b$. Dans
    $3(p+jq+j^{2}r)$, le coefficient de $a$ vaut
    $j(1-j)+j^{2}(1-j^{2})=j-j^{4}=0$, celui de $b$ vaut
    $(1-j^{2})+j^{2}(1-j)=1-j^{3}=0$, celui de $c$ vaut
    $(1-j)+j(1-j^{2})=1-j^{3}=0$. Donc $\res{p+jq+j^{2}r=0}$ : par la
    partie A, \res{PQR\text{ est équilatéral}} (et direct).
  \item $3(p+q+r)=\bigl((1-j^{2})+(1-j)\bigr)(a+b+c)
    =(2-j-j^{2})(a+b+c)=3(a+b+c)$, puisque $j+j^{2}=-1$ : les deux
    triangles ont le \res{\text{même centre de gravité}}.
\end{questions}
\end{corrige}
\end{blocLivre}

\begin{blocLivre}{cor-3-25b}
\begin{corrigeSuite}
\partieExo{C}{Longueurs et point de Fermat}
\begin{questions}[start=6]
  \item $j(a'-a)=-j^{3}b-j^{2}c-ja=-b-j^{2}c-ja=b'-b$, et de même
    $j(b'-b)=c'-c$. Multiplier par $j$, c'est tourner de
    $\frac{2\pi}{3}$ : $\res{AA'=BB'=CC'}$, et les vecteurs
    $\overrightarrow{AA'}$, $\overrightarrow{BB'}$, $\overrightarrow{CC'}$
    se déduisent l'un de l'autre par des rotations de $\frac{2\pi}{3}$ :
    les droites font deux à deux des angles de $\frac{\pi}{3}$.
  \item $P$ et $Q$ sont à distances $CP=\dfrac{\alpha}{\sqrt3}$ et
    $CQ=\dfrac{\beta}{\sqrt3}$ de $C$, et
    $\widehat{PCQ}=\dfrac{\pi}{6}+\widehat{C}+\dfrac{\pi}{6}$ (figure).
    Al-Kashi donne
    \[
    PQ^{2}=\frac{\alpha^{2}+\beta^{2}}{3}
    -\frac{2\alpha\beta}{3}\cos\Bigl(\widehat{C}+\frac{\pi}{3}\Bigr)
    =\frac{\alpha^{2}+\beta^{2}}{3}-\frac{\alpha\beta\cos\widehat{C}}{3}
    +\frac{\alpha\beta\sin\widehat{C}}{\sqrt3} .
    \]
    Avec $\alpha\beta\cos\widehat{C}=\frac{\alpha^{2}+\beta^{2}-\gamma^{2}}{2}$
    (Al-Kashi dans $ABC$) et $\alpha\beta\sin\widehat{C}=2\mathcal{S}$ :
    $\res{PQ^{2}=\dfrac{\alpha^{2}+\beta^{2}+\gamma^{2}}{6}
    +\dfrac{2}{\sqrt3}\,\mathcal{S}}$ — expression symétrique en
    $\alpha,\beta,\gamma$, ce qui redonne $PQ=QR=RP$.
  \item Pour $ABC$ équilatéral de côté $\ell$ :
    $\mathcal{S}=\frac{\sqrt3}{4}\ell^{2}$, donc
    $PQ^{2}=\frac{\ell^{2}}{2}+\frac{\ell^{2}}{2}$ et $\res{PQ=\ell}$.
    Chaque centre, par exemple $P$, est le symétrique de $A$ par rapport
    au centre $G$ de $ABC$ : $PQR$ est l'image de $ABC$ par la symétrie
    de centre $G$.
\end{questions}
\begin{center}
\begin{tikzpicture}[scale=0.95,line cap=round,line join=round]
  \coordinate (C) at (0,0);
  \coordinate (B) at (3.2,0);
  \coordinate (A) at (70:2.6);
  \coordinate (P) at (-30:1.848);
  \coordinate (Q) at (100:1.501);
  \draw[otInk,line width=0.9pt] (A) -- (C) -- (B);
  \draw[otMuted,line width=0.5pt,dash pattern=on 2pt off 1.5pt] (A) -- (B);
  \draw[otGreenDark,line width=1pt] (C) -- (P) (C) -- (Q);
  \draw[otGreenDark,line width=1.2pt] (P) -- (Q);
  \draw[otNavy,line width=0.6pt] (0.75,0) arc[start angle=0,end angle=70,radius=0.75];
  \draw[otGreen!70!black,line width=0.6pt] (0.55,0) arc[start angle=0,end angle=-30,radius=0.55];
  \draw[otGreen!70!black,line width=0.6pt] (70:0.55) arc[start angle=70,end angle=100,radius=0.55];
  \foreach \p in {A,B,C}{\fill[otInk] (\p) circle (1.4pt);}
  \foreach \p in {P,Q}{\fill[otGreenDark] (\p) circle (1.5pt);}
  \node[font=\small,anchor=north east] at (C) {$C$};
  \node[font=\small,anchor=west] at (B) {$B$};
  \node[font=\small,anchor=south west] at (A) {$A$};
  \node[font=\small,text=otGreenDark,anchor=north west] at (P) {$P$};
  \node[font=\small,text=otGreenDark,anchor=east] at (Q) {$Q$};
  \node[font=\scriptsize,text=otNavy] at (35:1.05) {$\widehat{C}$};
  \node[font=\scriptsize,text=otGreen!60!black] at (-15:0.85) {$\frac{\pi}{6}$};
  \node[font=\scriptsize,text=otGreen!60!black] at (85:0.85) {$\frac{\pi}{6}$};
\end{tikzpicture}\\[0.5mm]
\legende{Dans le triangle $PCQ$ : $CP=\alpha/\sqrt3$, $CQ=\beta/\sqrt3$ et
l'angle en $C$ vaut $\widehat{C}+\pi/3$.}
\end{center}
\end{corrigeSuite}
\end{blocLivre}
